\section{The Hodge metric on the gluing locus}\label{sec:moduli}

Return to the component-trivial bundles of
Section~\ref{sec:geometry}, and restrict the
quartic family to the one-parameter slice
$\tau_e=t$ for every edge. Thus
$L=\log(1/|t|)$ is common to all six necks.
We compare the actual harmonic-one-form
$L^2$ metric on a locus of flat-bundle moduli
with a finite quotient metric. The mass
weights in \eqref{eq:finite-forms} do not
enter this one-form comparison.

\subsection{The graph quotient and its embedding}
Write $\Gamma=K_4$, $d=r^2-1$ and
$g_\Gamma=b_1(\Gamma)=3$.
The irreducible graph representation space is
\[
 \mathcal M_\Gamma^{\mathrm{irr}}
   =\SU(r)^{E,\mathrm{irr}}/\SU(r)^V,
 \qquad U_{ij}\longmapsto g_iU_{ij}g_j^{-1}.
\]
After dividing the vertex gauge group by
its ineffective diagonal center, this is
a free proper action. The quotient is a
smooth real manifold of dimension
$(|E|-|V|)d=(g_\Gamma-1)d=2d$.
Choose an orientation for each edge.
A tangent representative is
$a_e=\dot U_eU_e^{-1}\in\mathfrak{su}(r)$.
Put
\begin{equation}\label{eq:graph-moduli-metric}
 (\delta_Ux)_{ij}=x_i-\Ad(U_{ij})x_j,\qquad
 G_{\Gamma,U}([a])
 =\min_{x\in\mathfrak{su}(r)^V}
                  \sum_e\|a_e+(\delta_Ux)_e\|^2 .
\end{equation}
Here $\langle X,Y\rangle=-\tr(XY)$.
This is the quotient metric of the
bi-invariant product metric on the edges.
The notation denotes its squared norm on
a tangent class.

Let $\mathcal M_t^{\mathrm{irr}}$ be the
moduli space of irreducible flat $\SU(r)$
connections on $C_t$. Its tangent space at
$D$ is $H^1(C_t;\operatorname{ad}D)$.
Irreducibility gives $H^0=0$ and duality
gives $H^2=0$, so the local quotient is
smooth. We use the Hodge metric: a class
has the squared $L^2$ norm of its unique
$D$-harmonic one-form representative,
equivalently the minimum of $\int|b|^2$
over its closed representatives modulo
exact forms. These deformation and
pairing conventions are those of
\cite[Section~1]{goldman}.
Via \cite{narasimhan-seshadri} this is
a metric on stable rank-$r$ degree-zero
bundles with trivial determinant.

Proposition~\ref{prop:bundle} gives a map
$\Phi_t:\mathcal M_\Gamma^{\mathrm{irr}}
\to\mathcal M_t^{\mathrm{irr}}$ by inflation
along $\pi_1(C_t)\twoheadrightarrow\pi_1(\Gamma)$.
It is an embedding. Surjectivity of this
group map shows that conjugate inflated
representations were already conjugate.
Evaluation on fixed lifts of the free
graph generators gives a continuous left
inverse on the image, smooth in local
representation slices. The differential
is injective: an inflated cocycle that
is a coboundary on the surface has the
same coboundary expression on these lifts.
This proves both immersion and the
topology of the embedding.
Denote the pullback Hodge metric by $G_t$.
The construction is a unitary Schottky
locus; the free-group construction itself
is classical \cite[Section~3, Proposition~1]{florentino}.

\begin{theorem}\label{thm:moduli-metric}
There is a constant $C$, depending only
on the fixed family, rank and exterior
collars, such that for large $L$,
\begin{equation}\label{eq:moduli-metric}
 \left(1-\frac CL\right)\frac{2\pi}{L}
       G_{\Gamma,U}([a])
 \le G_{t,U}([a])
 \le
 \left(1+\frac CL\right)\frac{2\pi}{L}
       G_{\Gamma,U}([a]).
\end{equation}
This holds for every irreducible $U$
and tangent class, uniformly as $U$
approaches reducibility.
At a reducible $U$ the same energy
inequalities hold on
$C^1(\Gamma;\operatorname{ad}U)/
\operatorname{im}\delta_U$; no smooth
moduli interpretation there is asserted.

For a connected relatively compact open
subset $\Omega\subset\mathcal M_\Gamma^{\mathrm{irr}}$,
let $d_{t,\Omega}$ and $d_{\Gamma,\Omega}$
be the intrinsic path distances using
paths in $\Omega$. Then
\begin{equation}\label{eq:moduli-distance}
 \sqrt{1-C/L}\,d_{\Gamma,\Omega}
 \le\sqrt{\frac{L}{2\pi}}\,d_{t,\Omega}
 \le\sqrt{1+C/L}\,d_{\Gamma,\Omega}.
\end{equation}
The image of $\Phi_t$ is Lagrangian in
the surface moduli space with its
cohomological symplectic form.
\end{theorem}

\subsection{Affine primitives and the energy infimum}
\begin{proof}
Fix a small collar radius $a$, independent
of $t$. The thick components and the
annuli of logarithmic length
$\ell_a=L+2\log a$ partition $C_t$.
Each has a flat unitary frame.

An inflated graph tangent class restricts
to zero cohomology on every thick piece
and every annulus: the represented family
has identically trivial holonomy on the
fundamental group of each such patch.
Thus any closed representative $b$ has
single-valued local primitives. In
the fixed flat frames write $b_i=df_i$.
On the left overlap choose the node
primitive to satisfy $f_e=f_i$.
On the right overlap it has the affine
jump
\begin{equation}\label{eq:affine-primitives}
 f_e=\Ad(U_{ij})f_j+a_e .
\end{equation}
The constants represent the graph
cochain class $[a]$. To check the sign,
if varying local parallel frames obey
$P_j(s)=P_i(s)U_{ij}(s)$, their
infinitesimal changes satisfy
$dk_i=-b_i$ and
$k_i=\Ad(U_{ij})k_j-a_{ij}$.
Taking $f_i=-k_i$ gives
\eqref{eq:affine-primitives}.
Changes of primitive constants change
the cochain by an element of
$\operatorname{im}\delta_U$.
Conversely these affine jumps construct
every graph cochain class.

Choose a minimizer $x$ in
\eqref{eq:graph-moduli-metric}.
Use constant component potentials
$-x_i$, and interpolate in log radius
on edge $ij$ from $-x_i$ to
$a_e-\Ad(U_{ij})x_j$.
The difference of the endpoints is
$a_e+(\delta_Ux)_e$.
Component derivatives vanish and
one-form energy is conformally invariant
in two dimensions. The resulting
energy is exactly
\[
 \frac{2\pi}{\ell_a}
        \sum_e\|a_e+(\delta_Ux)_e\|^2 .
\]
Piecewise affine primitives lie in
$H^1$ and give weakly closed $L^2$
one-forms. Smoothing the profiles
near their ends preserves the affine
jumps and approximates the same energy.
Equivalently one may use the weak
Hodge projection. The variational
characterization of $G_t$ gives the
upper estimate in
\eqref{eq:moduli-metric}.

The potentials $x_i$ themselves can be
large near a stabilizer change. Their
size does not enter this bound:
their component derivatives are zero
and the annular derivatives depend
only on the minimizing residual.

For the lower estimate take any smooth
closed representative $b$ of the class.
Let $x_i$ be the thick-component mean
of $f_i$, and $c_{ie}$ its boundary
circle mean. Fixed-domain Poincaré
and trace estimates give
\[
 \|c_{ie}-x_i\|
       \le C_a\|b\|_{L^2(\text{component }i)}.
\]
The constants are uniform in $t,U$:
the fixed thick metrics converge
smoothly and the frames are unitary.
Circle averaging and Cauchy--Schwarz
along the log cylinder give
\[
 E_e(b)\ge\frac{2\pi}{\ell_a}
       \|a_e+\Ad(U_{ij})c_{je}-c_{ie}\|^2 .
\]
Writing $E_{\mathrm n}$ and
$E_{\mathrm c}$ for the disjoint
neck and component energies, we obtain
\[
 \begin{split}
 G_{\Gamma,U}([a])^{1/2}
 &\le\left(\sum_e\|a_e-(\delta_Ux)_e\|^2\right)^{1/2}\\
 &\le\sqrt{\ell_a/(2\pi)}\,E_{\mathrm n}^{1/2}
                                  +C_aE_{\mathrm c}^{1/2}\\
 &\le\sqrt{\ell_a/(2\pi)+C_a^2}
                            \left(\int_{C_t}|b|^2\right)^{1/2}.
 \end{split}
\]
The last line uses $E_{\mathrm n}+E_{\mathrm c}
=\int|b|^2$. Taking the infimum over
representatives gives
\[
 G_{t,U}([a])
 \ge\left(\frac{\ell_a}{2\pi}+C_a^2\right)^{-1}
                                        G_{\Gamma,U}([a]),
\]
which is the required lower estimate.
One fixed collar radius suffices.
There is no division by an eigenvalue
of $\delta_U^*\delta_U$, and no
zero-form mass estimate is needed.
Fixed-fibre compact Hodge theory
identifies the infimum with the actual
harmonic norm. The same proof at a
reducible parameter is an inequality
on cohomology, independent of whether
the moduli quotient is smooth.

Taking square roots of the pointwise
quadratic inequalities and integrating
along a path proves the corresponding
length bounds. Taking infima over the
same paths in $\Omega$ proves
\eqref{eq:moduli-distance}.

Finally each graph class has a smooth
representative supported in the node
annuli with only a radial component.
Two such representatives wedge to
zero on each annulus; different
annuli are disjoint. The form
$\int-\tr(b_1\wedge b_2)$ therefore
vanishes on the image of $\Phi_t$.
All normalization components are
rational, so $g(C_t)=g_\Gamma=3$.
The image dimension $2d$ is half
the surface-moduli dimension $4d$,
proving the Lagrangian assertion.
\end{proof}

The Lagrangian interpretation concerns
classical Schottky geometry. The
quantitative assertion here is the
relative comparison of the actual
unweighted Hodge norm with the edge
quotient norm and its coefficient
$2\pi/L$.
Daskalopoulos--Wentworth
\cite[Proposition~3.23 and Corollary~3.24]{dw-local}
prove weighted operator continuity and
kernel continuation in a degeneration,
and obtain local moduli geometry.
Their operator analysis allows
arbitrary rank and retains boundary
zero modes. It does not state
\eqref{eq:moduli-metric}; the argument
above compares unweighted energy
infima directly. In particular it
does not invoke the limiting
zero-cokernel hypothesis used for
their kernel continuation.

The leading metric has no area weights
because it minimizes one-form energy.
The areas remain essential to the
zero-form spectrum and finite HYM
evolution. Equation~\eqref{eq:moduli-distance}
uses paths inside the indicated
locus; it makes no assertion about
shortcuts through the ambient
moduli space or distances across
reducible strata.
